% frontmatter/english/verano.tex -- instrucciones para el adulto del
% cuaderno de verano de "Read and Draw" (ediciones/english-verano.tex, ver
% notes/07-ediciones.md), en inglés como todo el cuaderno: van en vez de
% frontmatter/english/como-usar.tex, que habla del año entero. Detrás,
% las palabras de las instrucciones y "Who's who?" (los del libro entero)
% y, en vez del mapa del curso, el mapa del verano (\mapaEdicion,
% preamble.tex).
\section*{How to use this summer book}

This book is the summer of \textit{Read and Draw}: its last 13 weeks,
as a book of its own, for a little English reading every day of the
holidays. It is for a child who can already read in Spanish and has
learnt some English at school. There are 65 pages, one for every
weekday of the summer, and every page has the same three parts:

\begin{itemize}[leftmargin=6mm, itemsep=1mm]
\item \textbf{Date}, \textbf{Read} and \textbf{Notes}, for the grown-up.
  Tick how the reading went: alone, with help, or with difficulty.
\item \textbf{The text of the day}, in the green box: about 70 words.
  The sentences are short, but many of them have two parts, joined by
  \textit{because}, \textit{when}, \textit{who}, \textit{that},
  \textit{but} or \textit{so}, and the story is told in the past
  (\textit{was}, \textit{went}, \textit{saw}). The title of the box says
  how to read it.
\item \textbf{An activity about the text}, in the box below. Almost every
  day there is something to draw, colour, tick, circle or match, and to
  do it the child has to go back to the text: how many, what colour,
  where, why, what happened first. The golden rule is that \textbf{the
  answer is always in the text}. If the child can't find it, don't give
  it away: ask ``Where does it say that?'' and read the text again
  together.
\end{itemize}

\textbf{The activities.} \textbf{Read and draw}: draw exactly what the
text says -- a beautiful drawing is not the point, a true one is -- and
then check it with the questions at the bottom of the box.
\textbf{Read and colour}: colour a picture the way the text says.
\textbf{Where is it?}: draw things in the right place (\textit{in},
\textit{on}, \textit{under}, \textit{next to}...). \textbf{Yes or no?},
\textbf{Read and circle} and \textbf{Read and match}: short questions
about the details. \textbf{Match the halves}: join the two parts of a
long sentence from the text. \textbf{What happened first?}: put the
story in order. \textbf{The missing words}: find a sentence in the text
and write the word that is missing. \textbf{Word hunt}, \textbf{Riddle},
\textbf{Draw the story} (three pictures: the beginning, the middle and the
end of the week) and \textbf{Look back}. The next page shows the words
and the pictures of every instruction: read it together on the first
day.

\textbf{The story.} The pages tell the summer of \textbf{Lucía}, her
little brother \textbf{Dani} and their dog \textbf{Toby} in Grandma
Rosa's village -- and of their neighbours from London, \textbf{Amy},
her little brother \textbf{Sam} and \textbf{Pip}, their parrot who
talks. Amy spends July in London, so there are postcards, letters and
a diary; then the river, Andrés's bees, shooting stars, a storm and the
village festival. Then the Browns come to the village too, and Lucía
and Amy plan a treasure hunt -- and at the end it is back to the city,
and back to school.
Every week is a new chapter, and the page after ``Who's who?'' shows
them all, with a sun to colour for every day of reading.

\textbf{Ten minutes a day.} Read the text together first if it looks
long, and then let the child read it alone. When a word is new, explain
it with a gesture or a drawing rather than a translation, if you can.
Reading a page twice is always fine; so is doing the activity the next
morning. There is a diploma at the end of the book, and an answer key at
the back -- for the grown-up, to check, or to give a small clue.
\newpage
